\chapter{Kết luận và hướng phát triển}

Chương cuối cùng tổng kết những nội dung đã thực hiện trong dự án và chỉ ra các hướng cải tiến có thể triển khai trong tương lai. Phần kết luận tập trung vào mức độ hoàn thiện của hệ thống hiện tại, còn phần hướng phát triển nêu các cải tiến về thuật toán, đặc trưng và hạ tầng phục vụ.

\section{Kết luận}
Trong bài tập lớn này, đề tài đã xây dựng một hệ thống gợi ý sản phẩm thương mại điện tử theo kiến trúc hai giai đoạn. Hệ thống bắt đầu từ dữ liệu hành vi người dùng, thực hiện các bước tiền xử lý, chia phiên, gán nhãn, tạo đặc trưng, huấn luyện mô hình và cuối cùng triển khai API để trả về danh sách sản phẩm gợi ý. Qua đó, dự án thể hiện được một luồng xử lý tương đối đầy đủ của một hệ thống gợi ý, từ dữ liệu thô đến kết quả phục vụ người dùng.
\medskip

Ở giai đoạn xử lý dữ liệu, hệ thống đã chuyển đổi dữ liệu hành vi thành tập dữ liệu có cấu trúc phù hợp cho học máy. Các hành vi như xem sản phẩm, thêm vào giỏ hàng và mua hàng được ánh xạ thành nhãn và trọng số mẫu, giúp mô hình phân biệt được mức độ quan trọng của từng loại tương tác. Các đặc trưng người dùng, sản phẩm và phiên được xây dựng theo hướng point-in-time nhằm hạn chế rò rỉ dữ liệu tương lai và làm cho quá trình huấn luyện gần hơn với điều kiện thực tế.
\medskip

Ở giai đoạn mô hình, hệ thống sử dụng Matrix Factorization bằng PyTorch để học biểu diễn vector cho người dùng và sản phẩm. Các item embedding sau huấn luyện được đưa vào FAISS để hỗ trợ truy hồi ứng viên nhanh. Trên tập ứng viên này, mô hình LightGBM được sử dụng để chấm điểm và xếp hạng lại sản phẩm dựa trên nhiều nhóm đặc trưng khác nhau. Cách kết hợp Recall và Ranking giúp hệ thống cân bằng giữa tốc độ truy hồi và chất lượng danh sách gợi ý.
\medskip

Ở giai đoạn phục vụ, dự án triển khai API bằng FastAPI. Pipeline phục vụ suy luận nạp sẵn kho đặc trưng, mô hình Recall, FAISS index và mô hình Ranking vào bộ nhớ khi server khởi động. Khi nhận request, hệ thống có thể gợi ý theo sở thích dài hạn của người dùng và bổ sung thêm thông tin phiên hiện tại thông qua cơ chế Dual-Channel Recall. Ngoài ra, hệ thống cũng có cơ chế dự phòng cho người dùng mới và cơ chế tự căn chỉnh đặc trưng đầu vào cho LightGBM, giúp API ổn định hơn trong quá trình suy luận.
\medskip

Theo các kết quả thực nghiệm được ghi nhận trong tài liệu của dự án, hướng tiếp cận hai giai đoạn là phù hợp với bài toán gợi ý sản phẩm. Mô hình Recall có khả năng lọc nhanh các sản phẩm ứng viên, mô hình Ranking cải thiện thứ tự hiển thị, còn API đạt độ trễ dưới 50 ms trong các kịch bản thử nghiệm cục bộ đã ghi nhận. Dù vậy, hệ thống vẫn còn một số hạn chế liên quan đến kho đặc trưng, cập nhật dữ liệu thời gian thực và phương pháp đánh giá. Đây là cơ sở để tiếp tục cải tiến trong các phiên bản sau.

\section{Hướng phát triển}
Trong các hướng phát triển tiếp theo, mô hình Recall có thể được cải tiến bằng cách sử dụng BPR Loss thay cho cách học pointwise hiện tại \cite{rendle2009bpr}. BPR Loss phù hợp hơn với bản chất của bài toán gợi ý vì trực tiếp tối ưu quan hệ thứ tự giữa sản phẩm tích cực và sản phẩm tiêu cực của cùng một người dùng. Việc áp dụng phương pháp này có thể giúp tập ứng viên Recall có chất lượng tốt hơn.
\medskip

Đối với mô hình Ranking, hệ thống có thể chuyển từ binary objective sang các objective tối ưu trực tiếp cho xếp hạng như LambdaRank hoặc LambdaMART \cite{burges2010lambdamart}. Những phương pháp này được thiết kế để cải thiện các chỉ số như NDCG, do đó phù hợp hơn với mục tiêu đưa sản phẩm liên quan lên các vị trí đầu danh sách. Ngoài ra, mô hình Ranking cũng có thể được bổ sung các đặc trưng tương tác chéo như mức độ yêu thích của người dùng với từng danh mục, thương hiệu hoặc phân khúc giá.
\medskip

Về mặt hạ tầng, kho đặc trưng hiện tại mới được triển khai đơn giản bằng các file Parquet nạp vào RAM \cite{apacheparquet2026docs}. Khi dữ liệu lớn hơn hoặc cần cập nhật thường xuyên hơn, hệ thống nên chuyển sang một giải pháp feature store chuyên dụng, chẳng hạn Feast, để đồng bộ tốt hơn giữa đặc trưng offline dùng cho huấn luyện và đặc trưng online dùng cho phục vụ suy luận \cite{feast2026docs}. Bên cạnh đó, FAISS index hiện sử dụng \texttt{IndexFlatIP}, phù hợp với quy mô thử nghiệm nhưng có thể cần thay bằng HNSW hoặc IVF khi số lượng sản phẩm tăng lên rất lớn.
\medskip

Cuối cùng, hệ thống cần được đánh giá trong môi trường gần thực tế hơn. Ngoài các chỉ số offline như AUC, NDCG và MRR, nên bổ sung đánh giá trực tuyến (online evaluation) hoặc A/B testing để đo lường tác động của hệ thống đối với hành vi người dùng thật. Đây là bước cần thiết nếu muốn đưa hệ thống từ mức bài tập lớn hoặc thử nghiệm cục bộ sang môi trường sản phẩm thực tế.
